%\vspace{-15mm}
\begin{IEEEbiography}[
  {\includegraphics[width=1in,height=1.25in,clip,keepaspectratio]{bibio/pic_max.pdf}}]
  {Max Spahn} received the B.Sc. degree in Mechanical Engineering at the RWTH
  Aachen in 2018. He received his M.Sc. in General Mechanical Engineering at
  the RWTH in 2019. He also received the degree 
  Diplôme d'ingénieur des Arts et Manufactures
  in 2020 as part of a double degree program.
 He is currently working towards his Ph.D. degree in robotics. Furthermore, he
  is part of AIRLab, the AI for Retail Lab in Delft. His research interests
  in geometric approaches to motion planning and trajectory generation. He is
  passionate about open-sourcing scientific code, as he believes it will improve
  scientific outcome in robotics.
\end{IEEEbiography}
\vspace{-15mm}
\begin{IEEEbiography}[{\includegraphics[width=1in,height=1.25in,clip,keepaspectratio]{bibio/pic_martijn.pdf}}]{Martijn Wisse}
  received the M.Sc. and Ph.D. degrees in mechanical engineering from the Delft
  University of Technology, Delft, The Netherlands, in 2000 and 2004,
  respectively. He is currently a Professor at the Delft University of Technology.
  His previous research focused on passive dynamic walking robots and passive
  stability in the field of robot manipulators. He worked on underactuated
  grasping, open-loop stable manipulator control, design of robotic systems, and
  the creation of startups in this field. His current research interests focus on
  the neuroscientific principle of active inference and its application and
  advancements in robotics. 
\end{IEEEbiography}
\vspace{-15mm}
\begin{IEEEbiography}[{\includegraphics[width=1in,height=1.25in,clip,keepaspectratio]{bibio/pic_javier.pdf}}]{Javier Alonso-Mora}
  is an Associate Professor at the Delft University of
  Technology, where he leads the Autonomous Multi-robots Lab. He received his
  Ph.D. degree in robotics on cooperative motion planning
  from ETH Zurich, in partnership with Disney Research
  Zurich, and he was a Postdoctoral Associate at the Computer Science and
  Artificial Intelligence Lab (CSAIL) of the Massachusetts Institute of
  Technology. His research focuses on navigation, motion planning, and control of
  autonomous mobile robots, with a special emphasis on multi-robot systems, mobile
  manipulation, on-demand transportation, and robots that interact with other
  robots and humans in dynamic and uncertain environments. He currently serves as
  associate editor for IEEE Transactions on Robotics and for Springer Autonomous
  Robots. He is the recipient of a talent scheme VENI award from the Netherlands
  Organization for Scientific Research (2017), the ICRA Best Paper Award on
  Multi-robot Systems (2019) and an ERC Starting Grant (2021).
\end{IEEEbiography}
\vspace{27em}

